% ===========================================================================
\section{欧拉函数的求和恒等式}\label{sec:phisum}

\begin{theorem}[Gauss 恒等式]\label{thm:gauss-sum}
对任意正整数 $m$，
\[
  \sum_{d\mid m}\varphi(d)=m.
\]
\end{theorem}

\begin{proof}
把 $\Z_m$ 中的元素按它与 $m$ 的最大公因数分组（各组两两不交）：
\[
  m=|\Z_m|
  =\sum_{d\mid m}\bigl|\{i\in\Z_m : \gcd(i,m)=d\}\bigr|.
\]
对固定的 $d\mid m$，映射 $i\mapsto i/d$ 给出双射
\[
  \{i\in\Z_m : \gcd(i,m)=d\}\;\longleftrightarrow\;\Z_{m/d}^{*},
\]
因为 $\gcd(i,m)=d$ 当且仅当 $i=dj$ 且 $\gcd(j,m/d)=1$。故这一组恰有
$\varphi(m/d)$ 个元素，于是
\[
  m=\sum_{d\mid m}\bigl|\Z_{m/d}^{*}\bigr|
   =\sum_{d\mid m}\varphi\Bigl(\frac{m}{d}\Bigr)
   =\sum_{d\mid m}\varphi(d),
\]
最后一步只是对求和指标做代换 $d\leftrightarrow m/d$（$d$ 跑遍 $m$ 的所有
因子时，$m/d$ 也跑遍所有因子）。
\end{proof}

\begin{remark}[欧拉函数的计算（补充）]
对素数幂，模 $p^k$ 下与 $p^k$ 不互素的数恰为 $p$ 的倍数，故
\[
  \varphi(p^k)=p^k-p^{k-1}\qquad(k\ge1).
\]
结合 CRT 给出的积性（第 \ref{sec:crt} 节末的注记），对
$n=\prod_{i=1}^{r}p_i^{k_i}$ 有
\[
  \varphi(n)=\prod_{i=1}^{r}\bigl(p_i^{k_i}-p_i^{k_i-1}\bigr)
  =n\prod_{i=1}^{r}\Bigl(1-\frac{1}{p_i}\Bigr).
\]
例如 $\varphi(36)=36\bigl(1-\tfrac12\bigr)\bigl(1-\tfrac13\bigr)=12$。
\end{remark}
